% =====================================================================
%  var-task-rep -- Minimal proposal skeleton
%  One grey box per section. Add \subsection{...} freely inside any
%  section; each new subsection can carry its own box: \ph{ID}{title}{notes}.
% =====================================================================
\documentclass[11pt,a4paper]{article}

\usepackage[utf8]{inputenc}
\usepackage[T1]{fontenc}
\usepackage[margin=2.5cm]{geometry}
\usepackage{amsmath,amssymb,amsthm}
\usepackage{booktabs}
\usepackage{fvextra}
\DefineVerbatimEnvironment{promptbox}{Verbatim}{breaklines,breakanywhere,fontsize=\footnotesize,frame=single,framesep=2mm}
\usepackage{xcolor}
\usepackage[most]{tcolorbox}
\usepackage[round,authoryear]{natbib}
\usepackage[colorlinks=true,linkcolor=blue!50!black,citecolor=green!40!black,urlcolor=blue!60!black]{hyperref}

% Placeholder box: \ph{ID}{What goes here}{Notes}
\definecolor{phframe}{gray}{0.55}
\definecolor{phback}{gray}{0.95}
\newcommand{\ph}[3]{%
  \hypertarget{ph:#1}{}%
  \begin{tcolorbox}[enhanced,breakable,colback=phback,colframe=phframe,
    boxrule=0.4pt,arc=1mm,left=2mm,right=2mm,top=1mm,bottom=1mm,
    fonttitle=\sffamily\small\bfseries,coltitle=black,colbacktitle=gray!25,
    title={[#1]\quad TO WRITE --- #2}]
  \small #3
  \end{tcolorbox}}

% Inline value to be decided
\newcommand{\tbd}[1]{\textcolor{red!70!black}{\textsf{[TBD: #1]}}}

% Link back to the status table, placed under each section heading
\newcommand{\backtostatus}{\noindent{\hfill\footnotesize\hyperref[sec:status]{$\uparrow$\,status}}\par}

% Status row: \strow{label}{name}{content}{status}
\newcommand{\strow}[4]{\hyperref[#1]{#2} & #3 & #4 \\}

\newtheorem{definition}{Definition}
\newtheorem{hypothesis}{Hypothesis}

\title{Lexical Task Heads beyond Simple Tasks:\\ Logit Lens and Jacobian Lens Readouts}
\author{Filippo Merlo}
\date{Draft --- \today}

\begin{document}
\maketitle

\begin{abstract}
\phantomsection\label{sec:abstract}
Lexical task heads are attention heads whose output, decoded into the vocabulary, names the task the prompt asks for \citep{yang2026shared}. They were identified with the per-head logit lens on simple single-step tasks, such as mapping a country to its capital. We ask whether the same method finds them for complex tasks defined by explicit rules \citep{fu2026rules}, which need several inference steps, and whether Head-JLens \citep{merlo2026headjlens}, which decodes each head's output after mapping it through a Jacobian of the remaining layers averaged over many texts, finds lexical task heads that the logit lens misses. In Llama-3.1-8B-Instruct, we identify lexical task heads under both readouts on 14 simple tasks and 5 complex tasks at their easy level, with rule-based and example-based prompts, and test each head set by activation patching.
\end{abstract}

\section{Introduction}\label{sec:intro}
A language model's accuracy on a task can change a great deal with the wording of the prompt that asks for it: on Llama-3.1-8B-Instruct, prompts for the same antonym items differ in accuracy by up to a factor of four \citep[Fig.~1]{yang2026shared}. \citet{yang2026shared} explain part of this variation with \emph{lexical task heads}, attention heads whose outputs, decoded into the vocabulary with the logit lens, name the task (``opposite'' for antonyms, ``cities'' for capitals). The heads are largely shared between instruction-based and example-based prompts, patching their mean activation into a prompt without task information raises task accuracy, and prompts that fail activate them less.

Two limits of that work motivate this project. First, its 17 tasks are simple single-step mappings (country to capital, word to antonym), and the authors note that ``it remains unclear if these shared lexical mechanisms hold for more abstract tasks or tasks that are less easily defined by simple vocabulary clusters'' \citep[Sec.~4]{yang2026shared}. \citet{fu2026rules} provide such tasks: five novel tasks defined by explicit rules (a card game, a board game, a composed arithmetic operator, noun-class agreement in an artificial language, and lexical-category inference), each at three difficulty levels and each given as rules, as examples or as both. Second, the logit lens decodes a head's output directly and leaves out the computation of the layers that follow it; \citet[Sec.~4]{yang2026shared} note that task representations ``encoded in a non-lexical, abstract format'' may not project cleanly to tokens. Head-JLens reads each head through a Jacobian map of the remaining layers, and on Llama-3.1-8B-Instruct it ranks annotated intermediate words, which the prompt never states (such as the bridge entity of a two-hop question), higher than the per-head logit lens at most layers on five of six evaluation sets \citep[Sec.~3]{merlo2026headjlens}.

We therefore ask whether \citeauthor{yang2026shared}'s method identifies lexical task heads for the complex tasks of \citet{fu2026rules}, and how the logit lens and Head-JLens compare as readouts for identifying them.

\paragraph{Contributions.}
\begin{enumerate}
\item A test of the lexical-task-head method on complex, rule-defined tasks, with a causal check of the identified heads.
\item A comparison of the per-head logit lens and Head-JLens for identifying lexical task heads, on simple and complex tasks.
\end{enumerate}

\section{Related work}\label{sec:background}
\paragraph{Task representations across prompting styles.} Function vectors are sums of the mean outputs of a few attention heads that, added to the residual stream, trigger a task in another context \citep{todd2024function}. \citet{davidson2025different} extract them from instructions as well as from demonstrations and find that the two prompting styles engage mostly different heads: of the 20 top heads in each model, 2 to 7 are shared. \citet{yang2026shared} identify a different kind of task representation, lexical task heads, whose outputs decode to words describing the task; 73\% of them are shared between instruction-based and example-based prompts on average, and they are largely disjoint from function-vector heads \citep[App.~I]{yang2026shared}. We test the lexical-task-head method on complex tasks.

\paragraph{Reading attention heads in the vocabulary.} The logit lens decodes an intermediate state with the model's final normalization and unembedding \citep{nostalgebraist2020logitlens}; applied to one head's output, it is the readout of \citet{yang2026shared}. The Jacobian lens of \citet{gurnee2026workspace} instead maps a residual-stream state through a Jacobian of the remaining layers, averaged over many texts. Head-JLens moves this map to the concatenated head outputs, so that each head is read through the layers that follow it; in Llama-3.1-8B-Instruct it ranks annotated intermediate words higher than the per-head logit lens at most layers on five of six evaluation sets, and heads selected with it are better sites for edits \citep{merlo2026headjlens}. We compare the two readouts on a different use: identifying task-descriptive heads.

\paragraph{Learning tasks from rules and from examples.} \citet{fu2026rules} compare learning from rules, from examples and from both on five novel tasks. Rules outperform examples overall, adding examples to rules gives no significant gain, and the advantage of rules is largest for tasks that need algebraic computation. Their study is behavioural; we use their tasks to study task representations.

\section{Research questions}\label{sec:rq}
\begin{description}
\item[RQ1] Does the method of \citet{yang2026shared} identify lexical task heads for the complex tasks of \citet{fu2026rules} under rule-based and example-based prompts, and do the identified heads have a causal effect on task accuracy?
\item[RQ2] How do the per-head logit lens and Head-JLens compare as readouts for identifying lexical task heads, on simple and on complex tasks?
\item[RQ3] Are the lexical task heads of a complex task shared between rule-based and example-based prompts, as \citet{yang2026shared} find for simple tasks?
\end{description}

\begin{hypothesis}[replication]\label{h:rep}
On the simple tasks, Logit-head reproduces \citet{yang2026shared}: lexical task heads are found for every task, and patching them into prompts without task information raises accuracy above the unpatched baseline. \emph{Falsified} if, on the same model, the method finds no heads or no accuracy gain for tasks where \citet{yang2026shared} report both.
\end{hypothesis}

\begin{hypothesis}[complex tasks]\label{h:complex}
Following the limitation stated by \citet[Sec.~4]{yang2026shared}, the complex tasks yield fewer lexical task heads per task and a smaller causal effect than the simple tasks under Logit-head. \emph{Falsified} if they yield as many heads with as large an effect.
\end{hypothesis}

\begin{hypothesis}[readouts]\label{h:readout}
Following \citet{merlo2026headjlens}, J-head identifies at least the lexical task heads that Logit-head identifies, and the heads found only by J-head have a causal effect above baseline. \emph{Falsified} if the J-head set misses heads of the Logit-head set, or if the heads found only by J-head have no causal effect.
\end{hypothesis}

\paragraph{Scope.} We do not explain behavioural variation across wordings (the failure analyses of \citealp[Sec.~2.4]{yang2026shared}), and we use one model.

\section{Framework}\label{sec:framework}
\paragraph{Tasks and prompts.} A task $t$ is given by a dataset $D_t$ of inputs and correct outputs. For the complex tasks we use the easy level of each (Sec.~\ref{sec:tasks}). Each prompt for $t$ has a style $s\in\{\text{rule},\text{example},\text{combined}\}$: its task information is a rule text, a set of demonstrations, or both, followed by the input $x$. $P_{t,s}$ is the set of prompts of style $s$ for $t$, and $P^{+}_{t,s}\subseteq P_{t,s}$ those the model answers correctly and straight away, with the answer as the first generated text (in reasoning mode, Sec.~\ref{sec:reasoning}: correctly after reasoning).

\paragraph{Heads and readouts.} Following \citet[Sec.~2]{merlo2026headjlens}, $z_{l,h}(p)$ is the output of head $h$ in layer $l$ at the last position of $p$, before the output projection, and $c_{l,h}(p)=W^{(h)}_{O,l}z_{l,h}(p)$ its write to the residual stream. With $q_v=g\odot w_v$ the gain-scaled unembedding row of token $v$, the two readouts score $v$ by
\[
\ell^{\mathrm{L}}_{l,h}(p)[v]=q_v^\top c_{l,h}(p)\quad\text{(Logit-head)},\qquad
\ell^{\mathrm{J}}_{l,h}(p)[v]=q_v^\top J_{l,h}\,z_{l,h}(p)\quad\text{(J-head)}.
\]
$T^{R}_{k,l,h}(p)$ is the set of the $k$ highest-scoring tokens under readout $R\in\{\mathrm{L},\mathrm{J}\}$, and $V_t$ the term list of task $t$.

\begin{definition}[Lexical task head]\label{def:lth}
Following \citet[Sec.~2.2.2]{yang2026shared}, head $(l,h)$ is a lexical task head for $t$ on prompt $p$ under readout $R$ if $|T^{R}_{k,l,h}(p)\cap V_t|\geq n$. It is a lexical task head for $t$ and style $s$ if this holds for at least a fraction $\rho$ of the prompts in $P^{+}_{t,s}$. We use $n=1$, $k=10$, $\rho=0.1$, as \citet{yang2026shared}. $H^{R}_{t,s}$ denotes the set of such heads.
\end{definition}

\section{Methods}\label{sec:methods}

\subsection{Model}\label{sec:model}
Llama-3.1-8B-Instruct: it is the main model of \citet{yang2026shared}, and the only model with released Head-JLens matrices.

\subsection{Tasks}\label{sec:tasks}
\paragraph{Simple tasks.} The 14 tasks used by both \citet{yang2026shared} and \citet{davidson2025different}, with the datasets of \citet{todd2024function}: antonym, synonym, singular--plural, present--past, country--capital, country--currency, landmark--country, park--country, product--company, next-item, prev-item, English--French, English--German, English--Spanish.
Rule-based prompts: the six instruction templates per task of \citet[App.~A.3.1]{yang2026shared}.
Example-based prompts: \citet[App.~A.3]{yang2026shared}.

\paragraph{Complex tasks.} The five tasks of \citet{fu2026rules} at their easy level: Set Game, Tapatan, Operator Function, Noun Class Agreement, Lexical Category Inference. Data from their public release: 300 test items per task, drawn with their sampler and seed (all 600 for Lexical Category Inference), the same items under every prompt.
Rule-based prompts: their rules-only prompts \citep[App.~F]{fu2026rules} and further wordings (Sec.~\ref{sec:trmethods}).
Example-based and combined prompts: their minimum-coverage demonstration sets \citep[Sec.~2.5]{fu2026rules}.

\paragraph{Why these tasks.} Lexical task heads and their term lists were defined on the simple tasks \citep{yang2026shared}, so the method applies unchanged and our results can be compared with theirs. \citet[Sec.~4]{yang2026shared} restrict their study to tasks with simple lexical descriptions. The complex tasks are novel, need several inference steps, and are given in rule, example and combined forms that specify the same task on the same test items, so the task is held fixed while the way it is communicated varies \citep[Sec.~1--2]{fu2026rules}.

\paragraph{Answer format.} As in the source studies: the templates of \citet{yang2026shared} for the simple tasks; the answer formats and generation limits of \citet[App.~A]{fu2026rules} for the complex tasks, with greedy decoding (\citeauthor{fu2026rules} sample at temperature 0.1), so that the prompt whose heads are read and the scored answer always match. This keeps accuracy comparable with the reported results. \citeauthor{fu2026rules}'s instruction-tuned models often explain before answering, so forcing a direct answer would be a different setting. Lexical task heads are read at the last prompt token, where the answer is produced only if the model answers straight away; prompts on which it explains first are therefore left out of $P^{+}_{t,s}$ (Sec.~\ref{sec:framework}), except for the tasks run in reasoning mode (Sec.~\ref{sec:reasoning}).

\subsection{Task-representation methods}\label{sec:trmethods}
\paragraph{Lexical task heads} \citep[Sec.~2.2.2]{yang2026shared}.
Each attention head's output at the last prompt token is decoded with the two readouts of Sec.~\ref{sec:framework}: Logit-head, the per-head logit lens used by \citet{yang2026shared} \citep{nostalgebraist2020logitlens}, and J-head (Head-JLens), which applies the head's columns $J_{l,h}$ of the per-layer Jacobian map fitted on WikiText-103 \citep[Sec.~2]{merlo2026headjlens}. Fitted matrices for Llama-3.1-8B-Instruct are released with \citet[Sec.~S4]{merlo2026headjlens}.
Under each readout, lexical task heads are identified with Definition~\ref{def:lth}.
Term lists: \citet[App.~A.2.1]{yang2026shared} for the simple tasks; for the complex tasks, see below.
Causal check, as in \citet[Sec.~2.3.2, App.~G]{yang2026shared}: the mean of $c_{l,h}$ over $P^{+}_{t,s}$, multiplied by a scaling factor $\alpha$, replaces the output of every head in $H^{R}_{t,s}$ at the last position of a prompt without task information, and accuracy is measured for $\alpha$ from $-10$ to $10$. The patched prompts are those of the first 100 items answered correctly with one shot (simple tasks) or under the original rules prompt (complex tasks), and the mean leaves these items out. The control patches mean activations from a different task at the same heads: the control pairs of \citet{yang2026shared} for the simple tasks, and for the complex tasks a task with the same answer format (Noun Class Agreement and Lexical Category Inference with each other, Set Game and Operator Function with each other, Tapatan with Set Game).

\paragraph{What we add to the complex tasks.}\label{sec:adapt} \citet{fu2026rules} provide the data, one rules-only prompt per task and difficulty level, and the example-based and combined prompts. To apply lexical task heads to them we add three things.
\begin{enumerate}
\item \emph{Rule wordings}: six per original rule prompt (App.~\ref{app:orig}), each stating exactly the same rules. They are generated with GPT-5.5 through the Azure OpenAI API from a paraphrase prompt that keeps placeholders, output labels and the answer-format instruction verbatim (App.~\ref{app:genprompts}), and each is checked by us against the original; App.~\ref{app:wordings} lists them. The lexical-task-head criterion for a prompt style counts heads over at least 10\% of that style's prompts; \citet{yang2026shared} use six per task, and so do we.
\item \emph{Term lists} of task-descriptive words, of the kind in \citet[App.~A.2.1]{yang2026shared}: single words that describe the task rather than its answers, with case and number variants. They are generated with GPT-5.5 through the Azure OpenAI API from the rule prompt (App.~\ref{app:genprompts}), as \citet[App.~O]{yang2026shared} generated term lists with LLMs, and checked by us. Each word is expanded to token ids as in \citet{yang2026shared}: as written, with a leading space, in upper, capitalised and lower case, and after a \texttt{.}, \texttt{\_} or \texttt{-}, keeping every variant that is a single token of the model's tokenizer. There is one list per original rule prompt (App.~\ref{app:orig}): Tapatan's levels share one prompt and change only board size, piece count and line length, whereas the other tasks add rules or, for Operator Function, change the function \citep[Table~1, App.~F]{fu2026rules}. Lists are task-specific, not item-specific: for Lexical Category Inference, whose category definitions change from item to item, the list contains only words that describe the task itself. App.~\ref{app:terms} lists them.
\item \emph{Prompts without task information} for the causal check: the input block of the rules-only prompt alone (the lines that introduce the input and the answer, such as \texttt{Sentence: \{sentence\}} and \texttt{Answer:}), without rules or examples, the counterpart of the zero-shot prompts of \citet[Sec.~2.3.2]{yang2026shared}.
\end{enumerate}

\paragraph{Reasoning mode.}\label{sec:reasoning} For the complex tasks on which E0 finds that the model does not answer straight away (Set Game, Operator Function and Tapatan), we follow the instruction-tuned setup of \citet[App.~A]{fu2026rules} and let the model reason: chat template, greedy decoding, and an output limit of 6,000 new tokens, their limit for Operator Function, for all of these tasks. Answers are extracted as in their released code: after whitespace is collapsed, the first \texttt{\textbackslash boxed\{\}}, read as three cards of the form \texttt{(animal | biome | \ldots)} for Set Game and as an integer for Operator Function. Their code has no parser for Tapatan; we take the first of \texttt{A win}, \texttt{B win}, \texttt{continue} in the output. An output without an extractable answer within the limit is wrong. Generation stops once this answer is complete, since nothing after it can change the score. $P^{+}_{t,s}$ then contains the prompts whose final answer is correct, whatever precedes it. Definition~\ref{def:lth} is applied at two positions: the last prompt token, as for the other tasks, where heads show the task as identified from the prompt; and the pre-answer position, the token after which the first answer token is written, found by one forward pass over the prompt and the model's own output, where heads show the task as represented when the answer is produced. At the pre-answer position, the copy share also checks the generated text. In the causal check, each head set's scaled mean activation, taken at the position where the set was found, is patched at the last prompt position of a prompt without task information and kept at every generated position, and the final answer is scored. Heads at the last prompt token are compared across all tasks; heads at the pre-answer position and the accuracies after reasoning are reported separately. To bound the cost of long generations, these tasks use a fixed sample of 100 items per task, style and wording, and the causal check patches 10 items with $\alpha\in\{0,1,4,10\}$.

\paragraph{Why these readouts.} Lexical task heads locate a task representation and give its content in words \citep{yang2026shared}. Logit-head reproduces their readout. J-head reads each head through the layers that follow it: in Llama-3.1-8B-Instruct it ranks unstated intermediate words higher than Logit-head at most layers on five of six evaluation sets, and its readouts are better guides for edits \citep[Secs.~3--4]{merlo2026headjlens}. A lexical task head's decoded task word is likewise not the answer, and \citet[Sec.~4]{yang2026shared} note that their logit-lens criterion may miss task representations that do not project cleanly to tokens.

\subsection{Measures}\label{sec:measures}
For each task $t$, style $s$ and readout $R$:
\begin{itemize}
\item \emph{Accuracy} of each prompt, with the evaluation of the source paper: that of \citet{yang2026shared} for the simple tasks and that of \citet{fu2026rules} for the complex tasks.
\item \emph{Head count} $|H^{R}_{t,s}|$ and the distribution of the heads over layers.
\item \emph{Overlap} between two head sets $A$ and $B$: $|A\cap B|/|A|$ and $|A\cap B|/|B|$. Compared pairs: $H^{\mathrm{L}}_{t,s}$ and $H^{\mathrm{J}}_{t,s}$ (RQ2); $H^{R}_{t,\text{rule}}$ and $H^{R}_{t,\text{example}}$ (RQ3).
\item \emph{Causal effect}: the highest accuracy over $\alpha$ minus the unpatched accuracy, for the head set and for the control.
\item \emph{Copy share}: the share of prompts on which a head's matched terms ($T^{R}_{k,l,h}(p)\cap V_t$) occur in the prompt text, averaged over the heads of $H^{R}_{t,s}$. It separates heads that may copy task words from the prompt from heads that do not.
\end{itemize}

\section{Experiments}\label{sec:experiments}
\subsection{E0: Behavioural baseline}\label{sec:e0}
\emph{Question}: on which tasks does the model perform above chance? \emph{Procedure}: accuracy of every wording of every style, on all tasks, and the share of prompts on which the model does not answer straight away. \emph{Output}: the pairs of task and style used in E1 and E2, those whose accuracy exceeds chance, and the prompts $P^{+}_{t,s}$. The share of prompts left out decides whether other ways of handling them are needed.

\subsection{E1: Simple tasks}\label{sec:e1}
\emph{Question}: H\ref{h:rep}, and RQ2 on the simple tasks. \emph{Procedure}: identify $H^{\mathrm{L}}_{t,s}$ and $H^{\mathrm{J}}_{t,s}$ for the 14 simple tasks, for rule-based and example-based prompts (Yang's six instruction templates; 1, 5, 10, 20 and 30 shots); compute the measures of Sec.~\ref{sec:measures}; run the causal check on prompts without task information. \emph{Decides}: whether Logit-head reproduces \citet{yang2026shared} on our setup, and how J-head differs from it where the method is known to work.

\subsection{E2: Complex tasks}\label{sec:e2}
\emph{Question}: RQ1, H\ref{h:complex}, RQ2 on the complex tasks, and RQ3. \emph{Procedure}: as in E1, for each complex task kept by E0, with the added wordings and term lists (Sec.~\ref{sec:adapt}); example-based and combined prompts from \citet{fu2026rules}. \emph{Decides}: whether lexical task heads are found and causally effective for complex tasks, under which readout, and whether they are shared with example-based prompts.

\section{Plan}\label{sec:plan}
\paragraph{Steps.}
\begin{enumerate}
\item Build the additions of Sec.~\ref{sec:adapt}: rule wordings (each checked to state the same rules as the original prompt), term lists, prompts without task information.
\item Run E0 and select tasks.
\item Run E1.
\item Run E2.
\end{enumerate}
Reused code and data: the lexical-task-head code of \citet{yang2026shared}; the Head-JLens code and matrices of \citet{merlo2026headjlens}; the data and generation code of \citet{fu2026rules}; the task datasets of \citet{todd2024function}.

\paragraph{Limitations.} One model, and one difficulty level of each complex task. The term lists are generated by an LLM and checked by hand, and the identified heads depend on them. Rule-based prompts contain task words, so a head may decode them by copying from the prompt; the copy share measures this but does not decide it. Under J-head, some sites rank tokens highly on unrelated prompts as well \citep[Sec.~3.6]{merlo2026headjlens}, so a larger J-head set is not by itself evidence of more task representation; the causal check is needed. Prompts on which the model explains before answering are left out, which may remove some tasks or styles from E1 and E2.

\paragraph{Deliverables.} Head sets, measures and causal-check curves for every task, style and readout.

\bibliographystyle{plainnat}
\bibliography{references}

\appendix

% Appendix: prompts and generated material for the complex tasks.
% Included from proposal-minimal.tex.

\section{Original rule prompts}\label{app:orig}
The rules-only prompts of \citet{fu2026rules}, from their released code (\url{https://github.com/tinlaboratory/rules-vs-examples}, \texttt{rules-mode-prompts/}, MIT licence). Text in curly braces is filled per item. Tapatan uses the final-board-state prompt, whose levels differ only in its parameters: board size 3, 5, 7 and line length 3, 4, 5 for Easy, Medium, Hard \citep[Table~1]{fu2026rules}; the version in \citet[App.~F]{fu2026rules} also states the maximum number of pieces per player and that at most one player has a winning line. The Medium and Hard prompts of Lexical Category Inference are identical within each logic variant.

\subsection*{Set Game}

\paragraph{Easy.}
\begin{promptbox}
You will be shown 9 cards on the board and have to select 3 cards that form a GAME-SET based on the following rules.

RULES:
1. Each card has two attributes: an animal and a biome.
2. There are 10 types of animals: armadillo, axolotl, bat, cat, fox, horse, rabbit, sheep, turtle, wolf.
3. There are 10 types of biomes: plain, savanna, desert, swamp, forest, jungle, taiga, hill, badland, tundra.
4. A GAME-SET is a set of three cards: For each attribute (animal, biome), the three cards must be either ALL the SAME or ALL DIFFERENT. e.g. if 2 of the cards have the same value, and 1 of them has a different value, the set is NOT valid.

Here is the board:

{board}

Based on the rules, tell me which three cards here constitute a GAME-SET. Type the cards in the following format:
 \boxed{First card: CARD1
Second card: CARD2
Third card: CARD3}
\end{promptbox}

\paragraph{Medium.}
\begin{promptbox}
You will be shown 9 cards on the board and have to select 3 cards that form a GAME-SET based on the following rules.

RULES:
1. Each card has three attributes: an animal, a biome, and a food.
2. There are 10 types of animals: armadillo, axolotl, bat, cat, fox, horse, rabbit, sheep, turtle, wolf.
3. There are 10 types of biomes: plain, savanna, desert, swamp, forest, jungle, taiga, hill, badland, tundra.
4. There are 10 types of food: apple, potato, bread, cake, carrot, mutton, beef, cookie, pie, melon.
5. A GAME-SET is a set of three cards: For each attribute (animal, biome, food), the three cards must be either ALL the SAME or ALL DIFFERENT. e.g. if 2 of the cards have the same value, and 1 of them has a different value, the set is NOT valid.

Here is the board:

{board}

Based on the rules, tell me which three cards here constitute a GAME-SET. Type the cards in the following format:
 \boxed{First card: CARD1
Second card: CARD2
Third card: CARD3}
\end{promptbox}

\paragraph{Hard.}
\begin{promptbox}
You will be shown 9 cards on the board and have to select 3 cards that form a GAME-SET based on the following rules.

RULES:
1. Each card has a number and three attributes: an animal, a biome, and a food.
2. Numbers range from 1 to 12.
3. There are 10 types of animals: armadillo, axolotl, bat, cat, fox, horse, rabbit, sheep, turtle, wolf.
4. There are 10 types of biomes: plain, savanna, desert, swamp, forest, jungle, taiga, hill, badland, tundra.
5. There are 10 types of food: apple, potato, bread, cake, carrot, mutton, beef, cookie, pie, melon.
6. A GAME-SET is a set of three cards: For each attribute (animal, biome, food), the three cards must be either ALL the SAME or ALL DIFFERENT. e.g. if 2 of the cards have the same value, and 1 of them has a different value, the set is NOT valid.
7. But only for the number, 2 of the cards should have the same number, and 1 of them should have a different number in order for the set to be valid.

Here is the board:

{board}

Based on the rules, tell me which three cards here constitute a GAME-SET. Type the cards in the following format:
 \boxed{First card: CARD1
Second card: CARD2
Third card: CARD3}
\end{promptbox}

\subsection*{Tapatan}

\paragraph{all levels.}
\begin{promptbox}
You will be learning a new board game. You will then be given one final board state from that game. Determine the current game outcome.

RULES:
1. The board is a {board_size}x{board_size} grid of points with coordinates (x,y), where 0 <= x,y < {board_size}.
2. A board cell contains A for player A's piece, B for player B's piece, or . for an empty point.
3. Orthogonal neighbors are always adjacent. Diagonal neighbors are adjacent only from even-parity points where (x+y) is even: (x,y) connects to (x±1,y±1) when in bounds.
4. A player wins if they have {line_length} pieces in a straight line (horizontal, vertical, or diagonal) and each consecutive pair of points in that line is adjacent under rule 3.
5. If A has a valid winning line, output A win. If B has a valid winning line, output B win. If neither player has a valid winning line, output continue.
6. You must output a label that is exactly one of: A win, B win, or continue.

BOARD:
{board_state}

LABEL:
\end{promptbox}

\subsection*{Operator Function}

\paragraph{Easy.}
\begin{promptbox}
You will be shown a problem that uses a newly defined operator function based on the following rules and need to calculate the answer for the problem.

RULES:
1. The new operator function f is composed of two arithmetic operators, and its inputs are integers.
2. Let x, y, z be inputs of the function.
3. To compute f(x,y,z), first multiply x and y.
4. Then, add z to the result from the previous operation.
Problem: \({problem}\)
Answer (place the result in \boxed{}):\boxed{}
\end{promptbox}

\paragraph{Medium.}
\begin{promptbox}
You will be shown a problem that uses a newly defined operator function based on the following rules and need to calculate the answer for the problem.

RULES:
1. The new operator function g is composed of three arithmetic operators, and its inputs are integers.
2. Let x, y, z, a be inputs of the function.
3. To compute g(x,y,z,a), first subtract y from x.
4. Second, multiply the result from the previous operation by z.
5. Last, adds a to the result of the previous operation.
Problem: \({problem}\)
Answer (place the result in \boxed{}):\boxed{}
\end{promptbox}

\paragraph{Hard.}
\begin{promptbox}
You will be shown a problem that uses a newly defined operator function based on the following rules and need to calculate the answer for the problem.

RULES:
1. The new operator function h is composed of four arithmetic operators, and its inputs are integers.
2. Let x, y, z, a, b be inputs of the function.
3. To compute h(x,y,z,a,b), first add x and y.
4. Second, multiply the result from the previous operation by z.
5. Third, subtract a from the result of the previous operation.
6. Last, take the remainder when the result is divided by b.
Problem: \({problem}\)
Answer (place the result in \boxed{}):\boxed{}
\end{promptbox}

\subsection*{Noun Class Agreement}

\paragraph{Easy.}
\begin{promptbox}
You will be given one sentence in a synthetic language. Decide whether the sentence is grammatical.

RULES:
1. There are 2 noun classes: A and B.
2. Each noun phrase has the form DET NOUN.
3. The class A determiner is ka, and the class B determiner is ti.
4. The class A nouns are tac, sular, fep, wug. The class B nouns are bim, noko, glarn, zesh.
5. Sentences follow the template DET NOUN VERB DET NOUN.
6. Verbs are plain roots like dax, miv, lorp and never change.
7. A sentence is grammatical if and only if, in each noun phrase, the determiner matches the noun's class.
8. Your answer must be exactly one word: Yes or No.

Sentence: {sentence}
Answer:
\end{promptbox}

\paragraph{Medium.}
\begin{promptbox}
You will be given one sentence in a synthetic language. Decide whether the sentence is grammatical.

RULES:
1. There are 4 noun classes: A, B, C, D.
2. A noun phrase has the form DET ADJ-SUFF NOUN.
3. The class A determiner is ka, the class B determiner is ti, the class C determiner is su, and the class D determiner is vo.
4. Adjectives have the form STEM-SUFFIX. The suffix must match the noun's class.
5. The class A adjective suffix is -en, the class B adjective suffix is -os, the class C adjective suffix is -im, and the class D adjective suffix is -at.
6. Example adjective tokens are glim-en, prun-os, fen-im, zay-at.
7. The class A nouns are tav, fap. The class B nouns are bim, glarn. The class C nouns are sular, wug. The class D nouns are noko, zesh.
8. Sentences follow the template DET ADJ-SUFF NOUN VERB DET ADJ-SUFF NOUN.
9. Verbs are plain roots like dax, miv, lorp and never change.
10. A sentence is grammatical if and only if, in each noun phrase, both the determiner and the adjective suffix match the noun's class.
11. Your answer must be exactly one word: Yes or No.

Sentence: {sentence}
Answer:
\end{promptbox}

\paragraph{Hard.}
\begin{promptbox}
You will be given one sentence in a synthetic language. Decide whether the sentence is grammatical.

RULES:
1. There are 6 noun classes: A, B, C, D, E, F.
2. A noun phrase has the form DET ADJ-SUFF NOUN.
3. The class A determiner is ka, the class B determiner is ti, the class C determiner is su, the class D determiner is vo, the class E determiner is ne, and the class F determiner is la.
4. Adjectives have the form STEM-SUFFIX. The suffix must match the noun's class.
5. The class A adjective suffix is -en, the class B adjective suffix is -os, the class C adjective suffix is -im, the class D adjective suffix is -at, the class E adjective suffix is -uk, and the class F adjective suffix is -esh.
6. The verb has the form PREFIX-ROOT-SUFFIX.
7. The prefix must match the class of the subject noun, and the suffix must match the class of the object noun.
8. The class A subject prefix is ge-, the class B subject prefix is du-, the class C subject prefix is ri-, the class D subject prefix is zo-, the class E subject prefix is pa-, and the class F subject prefix is li-.
9. The class A object suffix is -an, the class B object suffix is -eb, the class C object suffix is -ig, the class D object suffix is -ot, the class E object suffix is -ul, and the class F object suffix is -er.
10. Example verb tokens are ge-dax-eb, pa-miv-ig, du-lorp-er.
11. The class A nouns are tav, fap. The class B nouns are bim, glarn. The class C nouns are sular, wug. The class D nouns are noko, zesh. The class E nouns are tac, fep. The class F nouns are lurn, prax.
12. Sentences follow the template DET ADJ-SUFF NOUN PREFIX-ROOT-SUFFIX DET ADJ-SUFF NOUN.
13. Verb roots are plain roots like dax, miv, lorp. Only the verb's prefix and suffix change.
14. A sentence is grammatical if and only if the subject noun phrase's determiner and adjective suffix match the subject noun's class, the object noun phrase's determiner and adjective suffix match the object noun's class, the verb prefix matches the subject noun's class, and the verb suffix matches the object noun's class.
15. Your answer must be exactly one word: Yes or No.

Sentence: {sentence}
Answer:
\end{promptbox}

\subsection*{Lexical Category Inference}

\paragraph{Easy.}
\begin{promptbox}
You will be shown a list of four words. Decide whether this list is correct based on the following rules.

RULES:
1. There are exactly four categories: {prompt_labels}.
2. Each category has the meaning shown in the category definitions below.
3. A correct list contains exactly four comma-separated items.
4. In a correct list, the first item is a member of Category 1, the second item is a member of Category 2, and so on.
5. If a category says "at least one", an item may satisfy any listed description. If it says "all", an item must satisfy every listed description.
6. Your answer must be exactly one word: Yes or No.

CATEGORY DEFINITIONS:
{label_block}

List: {candidate_block}
Answer:
\end{promptbox}

\paragraph{Medium and Hard, OR.}
\begin{promptbox}
You will be shown a list of four words. Decide whether this list is correct based on the following rules.

RULES:
1. There are exactly four categories: {prompt_labels}.
2. Each category has the meaning shown in the category definitions below.
3. A correct list contains exactly four comma-separated items.
4. In a correct list, the first item is a member of Category 1, the second item is a member of Category 2, and so on.
5. If a category lists multiple descriptions, an item is a member of that category if it satisfies at least one listed description.
6. Your answer must be exactly one word: Yes or No.

CATEGORY DEFINITIONS:
{label_block}

List: {candidate_block}
Answer:
\end{promptbox}

\paragraph{Medium and Hard, AND.}
\begin{promptbox}
You will be shown a list of four words. Decide whether this list is correct based on the following rules.

RULES:
1. There are exactly four categories: {prompt_labels}.
2. Each category has the meaning shown in the category definitions below.
3. A correct list contains exactly four comma-separated items.
4. In a correct list, the first item is a member of Category 1, the second item is a member of Category 2, and so on.
5. If a category lists multiple descriptions, an item is a member of that category only if it satisfies all listed descriptions.
6. Your answer must be exactly one word: Yes or No.

CATEGORY DEFINITIONS:
{label_block}

List: {candidate_block}
Answer:
\end{promptbox}

\section{Generation prompts}\label{app:genprompts}
Both prompts are given to GPT-5.5 through the Azure OpenAI API; \texttt{\{rule\_prompt\}} is an original prompt of App.~\ref{app:orig}.

\subsection*{Rule wordings}
\begin{promptbox}
You will be given a prompt that describes the rules of a task for a language model. Write {N} paraphrases of it.

Requirements:
1. Each paraphrase must state exactly the same rules as the original: the same constraints, conditions, quantities, names, symbols and labels. Do not add, remove, weaken or strengthen any rule, and do not add examples, hints or explanations.
2. Change the wording and the sentence structure. Reorder rules only where their order does not matter; if a rule refers to another rule by its number, keep the reference correct.
3. Copy verbatim every placeholder in curly braces (e.g. {board}), the answer-format instruction, and every output label or format string (e.g. "Yes or No", "A win, B win, or continue", \boxed{...}).
4. Keep unchanged the final lines that introduce the input and the answer (e.g. "Sentence: {sentence}" and "Answer:").
5. The paraphrases must differ from the original and from each other.

Return a JSON list of {N} strings.

Original prompt:
<<<
{rule_prompt}
>>>
\end{promptbox}

\subsection*{Term lists}
\begin{promptbox}
You will be given the rules of a task for a language model. List words that describe the task itself: what the task asks the model to do, its operation or relation, and its key concepts. For example, for the task "give the antonym of a word": opposite, reverse, antonym, contrary.

Requirements:
1. Single words only.
2. Do not list possible answers or output labels, words that only name the items of a particular input (e.g. specific animals, nouns or numbers), or placeholders.
3. Include singular and plural forms and capitalised forms. You may include translations into other languages.
4. Return a JSON list of strings.

Task rules:
<<<
{rule_prompt}
>>>
\end{promptbox}

\section{Rule wordings}\label{app:wordings}
The six wordings of each original rules prompt of App.~\ref{app:orig} at the easy level, the level in the study, generated and checked as in Sec.~\ref{sec:adapt}. Text in curly braces is filled per item.

% Generated by scripts/appendix_tex.py from materials/complex/wordings.json.

\subsection*{Set Game, easy level}

\paragraph{Wording 1.}
\begin{promptbox}
You will see 9 cards on the board, and your task is to choose 3 cards that make a GAME-SET according to these rules.

RULES:
1. Every card has two attributes: an animal and a biome.
2. The 10 animal types are: armadillo, axolotl, bat, cat, fox, horse, rabbit, sheep, turtle, wolf.
3. The 10 biome types are: plain, savanna, desert, swamp, forest, jungle, taiga, hill, badland, tundra.
4. A GAME-SET consists of three cards: For each attribute (animal, biome), the values across the three cards must be either ALL the SAME or ALL DIFFERENT. For example, if 2 cards share the same value and 1 card has a different value, the set is NOT valid.

Here is the board:

{board}

Based on the rules, tell me which three cards here constitute a GAME-SET. Type the cards in the following format:
 \boxed{First card: CARD1
Second card: CARD2
Third card: CARD3}
\end{promptbox}

\paragraph{Wording 2.}
\begin{promptbox}
You are given a board with 9 cards and must pick 3 cards that constitute a GAME-SET under the rules below.

RULES:
1. Each card contains two attributes: an animal and a biome.
2. The animals come in 10 types: armadillo, axolotl, bat, cat, fox, horse, rabbit, sheep, turtle, wolf.
3. The biomes come in 10 types: plain, savanna, desert, swamp, forest, jungle, taiga, hill, badland, tundra.
4. A GAME-SET is a set of three cards: For each attribute (animal, biome), the three cards must have values that are either ALL the SAME or ALL DIFFERENT. For instance, when 2 cards have the same value and 1 card has another value, the set is NOT valid.

Here is the board:

{board}

Based on the rules, tell me which three cards here constitute a GAME-SET. Type the cards in the following format:
 \boxed{First card: CARD1
Second card: CARD2
Third card: CARD3}
\end{promptbox}

\paragraph{Wording 3.}
\begin{promptbox}
A board containing 9 cards will be shown, and you need to select 3 cards that form a GAME-SET by following the rules.

RULES:
1. Each card is defined by two attributes: an animal and a biome.
2. There are 10 animal types: armadillo, axolotl, bat, cat, fox, horse, rabbit, sheep, turtle, wolf.
3. There are 10 biome types: plain, savanna, desert, swamp, forest, jungle, taiga, hill, badland, tundra.
4. A GAME-SET is made up of three cards: For each attribute (animal, biome), the three cards must be either ALL the SAME or ALL DIFFERENT. If, for example, 2 of the cards have one value and 1 has a different value, the set is NOT valid.

Here is the board:

{board}

Based on the rules, tell me which three cards here constitute a GAME-SET. Type the cards in the following format:
 \boxed{First card: CARD1
Second card: CARD2
Third card: CARD3}
\end{promptbox}

\paragraph{Wording 4.}
\begin{promptbox}
You will be presented with 9 cards on the board and must identify 3 cards that create a GAME-SET according to the following rules.

RULES:
1. Each card has two attributes: an animal and a biome.
2. The 10 kinds of animals are: armadillo, axolotl, bat, cat, fox, horse, rabbit, sheep, turtle, wolf.
3. The 10 kinds of biomes are: plain, savanna, desert, swamp, forest, jungle, taiga, hill, badland, tundra.
4. A GAME-SET is a group of three cards: For each attribute (animal, biome), the three cards must be either ALL the SAME or ALL DIFFERENT. e.g. when 2 of the cards share the same value and 1 of them has a different value, the set is NOT valid.

Here is the board:

{board}

Based on the rules, tell me which three cards here constitute a GAME-SET. Type the cards in the following format:
 \boxed{First card: CARD1
Second card: CARD2
Third card: CARD3}
\end{promptbox}

\paragraph{Wording 5.}
\begin{promptbox}
From the 9 cards displayed on the board, choose 3 cards that satisfy the definition of a GAME-SET using these rules.

RULES:
1. Each card includes two attributes: an animal and a biome.
2. There are 10 types of animals: armadillo, axolotl, bat, cat, fox, horse, rabbit, sheep, turtle, wolf.
3. There are 10 types of biomes: plain, savanna, desert, swamp, forest, jungle, taiga, hill, badland, tundra.
4. A GAME-SET is a collection of three cards: For each attribute (animal, biome), the three cards must be either ALL the SAME or ALL DIFFERENT. e.g. if 2 of the cards have the same value while 1 of them has a different value, the set is NOT valid.

Here is the board:

{board}

Based on the rules, tell me which three cards here constitute a GAME-SET. Type the cards in the following format:
 \boxed{First card: CARD1
Second card: CARD2
Third card: CARD3}
\end{promptbox}

\paragraph{Wording 6.}
\begin{promptbox}
You will be shown a board of 9 cards, and you have to determine 3 cards that form a GAME-SET based on the rules below.

RULES:
1. Each card has two attributes: an animal and a biome.
2. The animal attribute has 10 possible types: armadillo, axolotl, bat, cat, fox, horse, rabbit, sheep, turtle, wolf.
3. The biome attribute has 10 possible types: plain, savanna, desert, swamp, forest, jungle, taiga, hill, badland, tundra.
4. A GAME-SET is a set containing three cards: For each attribute (animal, biome), the three cards must be either ALL the SAME or ALL DIFFERENT. e.g. if 2 of the cards have matching values and 1 of them has a different value, the set is NOT valid.

Here is the board:

{board}

Based on the rules, tell me which three cards here constitute a GAME-SET. Type the cards in the following format:
 \boxed{First card: CARD1
Second card: CARD2
Third card: CARD3}
\end{promptbox}

\subsection*{Operator Function, easy level}

\paragraph{Wording 1.}
\begin{promptbox}
A problem involving a newly defined operator function will be presented, and you must determine its answer using these rules.

RULES:
1. The new operator function f consists of two arithmetic operators, and all of its inputs are integers.
2. Let x, y, z be the inputs to the function.
3. When computing f(x,y,z), begin by multiplying x and y.
4. After that, add z to the outcome of the preceding operation.
Problem: \({problem}\)
Answer (place the result in \boxed{}):\boxed{}
\end{promptbox}

\paragraph{Wording 2.}
\begin{promptbox}
You need to solve a problem that uses a newly introduced operator function defined according to the rules below.

RULES:
1. The new operator function f is made from two arithmetic operators, and its inputs are integers.
2. Let x, y, z represent the inputs of the function.
3. To evaluate f(x,y,z), multiply x by y first.
4. Next, add z to the result obtained from the previous operation.
Problem: \({problem}\)
Answer (place the result in \boxed{}):\boxed{}
\end{promptbox}

\paragraph{Wording 3.}
\begin{promptbox}
A problem will be given that relies on a newly defined operator function; compute the answer by applying the following rules.

RULES:
1. The new operator function f is formed using two arithmetic operators, and the inputs to it are integers.
2. Let x, y, z be inputs of the function.
3. For f(x,y,z), first take the product of x and y.
4. Then add z to the result of that prior operation.
Problem: \({problem}\)
Answer (place the result in \boxed{}):\boxed{}
\end{promptbox}

\paragraph{Wording 4.}
\begin{promptbox}
You will receive a problem using a newly defined operator function, and you must calculate the result according to these rules.

RULES:
1. The new operator function f is built from two arithmetic operators, and its inputs are integers.
2. Let x, y, z denote the inputs of the function.
3. In computing f(x,y,z), the first step is to multiply x and y.
4. The next step is to add z to the result from the previous operation.
Problem: \({problem}\)
Answer (place the result in \boxed{}):\boxed{}
\end{promptbox}

\paragraph{Wording 5.}
\begin{promptbox}
The problem shown will involve a newly defined operator function, and the answer must be found using the rules that follow.

RULES:
1. The new operator function f is composed of two arithmetic operators, with integer inputs.
2. Let x, y, z be the function's inputs.
3. To calculate f(x,y,z), start by multiplying x and y.
4. Then add z to the value produced by the preceding operation.
Problem: \({problem}\)
Answer (place the result in \boxed{}):\boxed{}
\end{promptbox}

\paragraph{Wording 6.}
\begin{promptbox}
Use the rules below to calculate the answer to a problem involving a newly defined operator function.

RULES:
1. The new operator function f contains two arithmetic operators, and the inputs are integers.
2. Let x, y, z be inputs of the function.
3. To compute f(x,y,z), multiply x and y before anything else.
4. Then add z to the result from the previous operation.
Problem: \({problem}\)
Answer (place the result in \boxed{}):\boxed{}
\end{promptbox}

\subsection*{Noun Class Agreement, easy level}

\paragraph{Wording 1.}
\begin{promptbox}
You will receive a single sentence in a synthetic language. Determine whether it is grammatical.

RULES:
1. The language has 2 noun classes: A and B.
2. Every noun phrase is structured as DET NOUN.
3. ka is the determiner for class A, and ti is the determiner for class B.
4. The nouns in class A are tac, sular, fep, wug. The nouns in class B are bim, noko, glarn, zesh.
5. Sentences use the pattern DET NOUN VERB DET NOUN.
6. The verbs are plain roots like dax, miv, lorp and never change.
7. A sentence is grammatical exactly when, in every noun phrase, the determiner matches the noun's class.
8. Your answer must be exactly one word: Yes or No.

Sentence: {sentence}
Answer:
\end{promptbox}

\paragraph{Wording 2.}
\begin{promptbox}
You will be provided with one sentence written in a synthetic language. Judge whether the sentence is grammatical.

RULES:
1. There are 2 noun classes: A and B.
2. A noun phrase always has the form DET NOUN.
3. The determiner for class A is ka, and the determiner for class B is ti.
4. Class A contains the nouns tac, sular, fep, wug. Class B contains the nouns bim, noko, glarn, zesh.
5. A sentence follows the template DET NOUN VERB DET NOUN.
6. Verbs are plain roots like dax, miv, lorp and never change.
7. A sentence is grammatical if and only if the determiner in each noun phrase matches the noun's class.
8. Your answer must be exactly one word: Yes or No.

Sentence: {sentence}
Answer:
\end{promptbox}

\paragraph{Wording 3.}
\begin{promptbox}
You will see one sentence in a synthetic language. Decide if that sentence is grammatical.

RULES:
1. The noun classes are 2 in number: A and B.
2. Each noun phrase is of the form DET NOUN.
3. ka is used as the class A determiner, and ti is used as the class B determiner.
4. tac, sular, fep, wug are the class A nouns. bim, noko, glarn, zesh are the class B nouns.
5. Sentences have the template DET NOUN VERB DET NOUN.
6. Verbs are plain roots like dax, miv, lorp and never change.
7. A sentence is grammatical if and only if each noun phrase has a determiner that matches the noun's class.
8. Your answer must be exactly one word: Yes or No.

Sentence: {sentence}
Answer:
\end{promptbox}

\paragraph{Wording 4.}
\begin{promptbox}
A single sentence in a synthetic language will be given to you. Decide whether it is grammatical.

RULES:
1. There are 2 noun classes, A and B.
2. Noun phrases take the form DET NOUN.
3. The class A determiner is ka; the class B determiner is ti.
4. The class A nouns are tac, sular, fep, wug. The class B nouns are bim, noko, glarn, zesh.
5. Sentences are built according to DET NOUN VERB DET NOUN.
6. Verbs are plain roots like dax, miv, lorp and never change.
7. A sentence is grammatical if and only if, for each noun phrase, the determiner matches the noun's class.
8. Your answer must be exactly one word: Yes or No.

Sentence: {sentence}
Answer:
\end{promptbox}

\paragraph{Wording 5.}
\begin{promptbox}
You will get one sentence from a synthetic language. Determine whether the sentence is grammatical.

RULES:
1. There are 2 noun classes: A and B.
2. The form of every noun phrase is DET NOUN.
3. For class A, the determiner is ka, and for class B, the determiner is ti.
4. The class A nouns are tac, sular, fep, wug. The class B nouns are bim, noko, glarn, zesh.
5. Sentences conform to the template DET NOUN VERB DET NOUN.
6. Verbs are plain roots like dax, miv, lorp and never change.
7. A sentence is grammatical if and only if the determiner matches the noun's class in each noun phrase.
8. Your answer must be exactly one word: Yes or No.

Sentence: {sentence}
Answer:
\end{promptbox}

\paragraph{Wording 6.}
\begin{promptbox}
You will be shown one sentence in a synthetic language. Decide whether that sentence is grammatical.

RULES:
1. The language includes 2 noun classes: A and B.
2. Each noun phrase follows the form DET NOUN.
3. ka is the class A determiner, while ti is the class B determiner.
4. The class A nouns are tac, sular, fep, wug. The class B nouns are bim, noko, glarn, zesh.
5. Sentences follow this template: DET NOUN VERB DET NOUN.
6. Verbs are plain roots like dax, miv, lorp and never change.
7. A sentence is grammatical if and only if, within every noun phrase, the determiner matches the noun's class.
8. Your answer must be exactly one word: Yes or No.

Sentence: {sentence}
Answer:
\end{promptbox}

\subsection*{Tapatan, easy level}

\paragraph{Wording 1.}
\begin{promptbox}
You are going to learn a new board game. After that, you will receive a single final board state from the game. Decide the present game outcome.

RULES:
1. The board is a {board_size}x{board_size} grid of points whose coordinates are (x,y), with 0 <= x,y < {board_size}.
2. Each board cell holds A for player A's piece, B for player B's piece, or . for an empty point.
3. Orthogonal neighbors are always adjacent. Diagonal neighbors are adjacent only from even-parity points where (x+y) is even: (x,y) connects to (x±1,y±1) when in bounds.
4. A player wins when they have {line_length} pieces in one straight line (horizontal, vertical, or diagonal), with every consecutive pair of points in that line adjacent according to rule 3.
5. If A has a valid winning line, output A win. If B has a valid winning line, output B win. If neither player has a valid winning line, output continue.
6. You must output a label that is exactly one of: A win, B win, or continue.

BOARD:
{board_state}

LABEL:
\end{promptbox}

\paragraph{Wording 2.}
\begin{promptbox}
You will be taught a new board game, and then you will be shown one final board state from that game. Identify the game outcome at that moment.

RULES:
1. The board consists of a {board_size}x{board_size} grid of points at coordinates (x,y), where 0 <= x,y < {board_size}.
2. A board cell contains A for player A's piece, B for player B's piece, or . for an empty point.
3. Orthogonal neighbors are always adjacent. Diagonal neighbors are adjacent only from even-parity points where (x+y) is even: (x,y) connects to (x±1,y±1) when in bounds.
4. A player wins if they possess {line_length} pieces in a straight line (horizontal, vertical, or diagonal) and each successive pair of points in that line is adjacent under rule 3.
5. If A has a valid winning line, output A win. If B has a valid winning line, output B win. If neither player has a valid winning line, output continue.
6. You must output a label that is exactly one of: A win, B win, or continue.

BOARD:
{board_state}

LABEL:
\end{promptbox}

\paragraph{Wording 3.}
\begin{promptbox}
You will learn the rules for a new board game. You will then be provided with one final board state from it. Determine what the current outcome of the game is.

RULES:
1. The board is a grid of points with size {board_size}x{board_size} and coordinates (x,y), where 0 <= x,y < {board_size}.
2. In a board cell, A denotes player A's piece, B denotes player B's piece, and . denotes an empty point.
3. Orthogonal neighbors are always adjacent. Diagonal neighbors are adjacent only from even-parity points where (x+y) is even: (x,y) connects to (x±1,y±1) when in bounds.
4. A player wins by having {line_length} pieces in a straight line (horizontal, vertical, or diagonal), provided that each consecutive pair of points in that line is adjacent under rule 3.
5. If A has a valid winning line, output A win. If B has a valid winning line, output B win. If neither player has a valid winning line, output continue.
6. You must output a label that is exactly one of: A win, B win, or continue.

BOARD:
{board_state}

LABEL:
\end{promptbox}

\paragraph{Wording 4.}
\begin{promptbox}
You will be introduced to a new board game. Next, one final board state from the game will be given to you. Work out the current game outcome.

RULES:
1. The board is a {board_size}x{board_size} grid made of points with coordinates (x,y), where 0 <= x,y < {board_size}.
2. A board cell contains A for player A's piece, B for player B's piece, or . for an empty point.
3. Orthogonal neighbors are always adjacent. Diagonal neighbors are adjacent only from even-parity points where (x+y) is even: (x,y) connects to (x±1,y±1) when in bounds.
4. A player wins if {line_length} of their pieces form a straight line (horizontal, vertical, or diagonal) and every consecutive pair of points in that line is adjacent under rule 3.
5. If A has a valid winning line, output A win. If B has a valid winning line, output B win. If neither player has a valid winning line, output continue.
6. You must output a label that is exactly one of: A win, B win, or continue.

BOARD:
{board_state}

LABEL:
\end{promptbox}

\paragraph{Wording 5.}
\begin{promptbox}
You will receive the rules of a new board game. Then you will be given one final state of the board from that game. Determine the outcome currently shown.

RULES:
1. The board is a {board_size}x{board_size} point grid using coordinates (x,y), where 0 <= x,y < {board_size}.
2. A board cell contains A for player A's piece, B for player B's piece, or . for an empty point.
3. Orthogonal neighbors are always adjacent. Diagonal neighbors are adjacent only from even-parity points where (x+y) is even: (x,y) connects to (x±1,y±1) when in bounds.
4. A win occurs for a player if they have {line_length} pieces in a straight line (horizontal, vertical, or diagonal) and each consecutive pair of points in that line is adjacent under rule 3.
5. If A has a valid winning line, output A win. If B has a valid winning line, output B win. If neither player has a valid winning line, output continue.
6. You must output a label that is exactly one of: A win, B win, or continue.

BOARD:
{board_state}

LABEL:
\end{promptbox}

\paragraph{Wording 6.}
\begin{promptbox}
You will study a new board game and then receive one final board state from that game. Decide which current game outcome applies.

RULES:
1. The board is a {board_size}x{board_size} grid of points, each with coordinates (x,y), where 0 <= x,y < {board_size}.
2. A board cell contains A for player A's piece, B for player B's piece, or . for an empty point.
3. Orthogonal neighbors are always adjacent. Diagonal neighbors are adjacent only from even-parity points where (x+y) is even: (x,y) connects to (x±1,y±1) when in bounds.
4. A player wins if they have {line_length} pieces arranged in a straight line (horizontal, vertical, or diagonal) and each consecutive pair of points in that line is adjacent under rule 3.
5. If A has a valid winning line, output A win. If B has a valid winning line, output B win. If neither player has a valid winning line, output continue.
6. You must output a label that is exactly one of: A win, B win, or continue.

BOARD:
{board_state}

LABEL:
\end{promptbox}

\subsection*{Lexical Category Inference, easy level}

\paragraph{Wording 1.}
\begin{promptbox}
You will receive a list made up of four words. Determine whether the list is correct according to these rules.

RULES:
1. There are exactly four categories: {prompt_labels}.
2. The meaning of each category is given in the category definitions below.
3. A correct list has exactly four items separated by commas.
4. For a list to be correct, item one must belong to Category 1, item two must belong to Category 2, and this pattern continues in order.
5. When a category says "at least one", an item may match any of the listed descriptions. When it says "all", an item must match every listed description.
6. Your answer must be exactly one word: Yes or No.

CATEGORY DEFINITIONS:
{label_block}

List: {candidate_block}
Answer:
\end{promptbox}

\paragraph{Wording 2.}
\begin{promptbox}
A list of four words will be presented to you. Judge whether it is correct using the rules below.

RULES:
1. There are exactly four categories: {prompt_labels}.
2. Each category means what is specified in the category definitions below.
3. A correct list consists of exactly four comma-separated items.
4. In a correct list, the items correspond in order: the first item is in Category 1, the second item is in Category 2, and so on.
5. If a category says "at least one", an item may satisfy any listed description. If it says "all", an item must satisfy every listed description.
6. Your answer must be exactly one word: Yes or No.

CATEGORY DEFINITIONS:
{label_block}

List: {candidate_block}
Answer:
\end{promptbox}

\paragraph{Wording 3.}
\begin{promptbox}
You will be given four words as a list. Decide if that list is correct by applying the following rules.

RULES:
1. There are exactly four categories: {prompt_labels}.
2. The category definitions below state the meaning of each category.
3. To be correct, a list must include exactly four items separated by commas.
4. To be correct, the first item must be a member of Category 1, the second item must be a member of Category 2, and so forth.
5. If a category says "at least one", an item may satisfy any listed description. If it says "all", an item must satisfy every listed description.
6. Your answer must be exactly one word: Yes or No.

CATEGORY DEFINITIONS:
{label_block}

List: {candidate_block}
Answer:
\end{promptbox}

\paragraph{Wording 4.}
\begin{promptbox}
You will see a four-word list. Assess whether this list is correct under the rules that follow.

RULES:
1. There are exactly four categories: {prompt_labels}.
2. Each category has the meaning provided in the category definitions below.
3. A list is correct only if it contains exactly four comma-separated items.
4. A correct list places a Category 1 member first, a Category 2 member second, and continues in that order.
5. If a category says "at least one", an item may satisfy any listed description. If it says "all", an item must satisfy every listed description.
6. Your answer must be exactly one word: Yes or No.

CATEGORY DEFINITIONS:
{label_block}

List: {candidate_block}
Answer:
\end{promptbox}

\paragraph{Wording 5.}
\begin{promptbox}
You will be shown four words in list form. Decide whether the list qualifies as correct according to these rules.

RULES:
1. There are exactly four categories: {prompt_labels}.
2. The definitions below give the meaning for each category.
3. Exactly four comma-separated items are required for a correct list.
4. In a correct list, the first, second, and later items must respectively be members of Category 1, Category 2, and the following categories in sequence.
5. If a category says "at least one", an item may satisfy any listed description. If it says "all", an item must satisfy every listed description.
6. Your answer must be exactly one word: Yes or No.

CATEGORY DEFINITIONS:
{label_block}

List: {candidate_block}
Answer:
\end{promptbox}

\paragraph{Wording 6.}
\begin{promptbox}
A list containing four words will be provided. Determine whether it is correct based on the rules listed here.

RULES:
1. There are exactly four categories: {prompt_labels}.
2. Each category's meaning is the one shown in the category definitions below.
3. For a list to be correct, it must contain exactly four items, with commas separating them.
4. For a list to be correct, its first item must belong to Category 1, its second item must belong to Category 2, and the remaining items must follow the same order.
5. If a category says "at least one", an item may satisfy any listed description. If it says "all", an item must satisfy every listed description.
6. Your answer must be exactly one word: Yes or No.

CATEGORY DEFINITIONS:
{label_block}

List: {candidate_block}
Answer:
\end{promptbox}


\section{Term lists}\label{app:terms}
The term list of each complex task at the easy level, generated and checked as in Sec.~\ref{sec:adapt}, as words before their expansion to token ids. Tapatan's list is made from its single prompt, shared by all levels.

\begin{center}
% Generated by scripts/appendix_tex.py from materials/complex/terms.json.
\begin{footnotesize}
\begin{tabular}{@{}lp{0.7\textwidth}@{}}
\toprule
Task (easy level) & Terms \\
\midrule
Set Game & set, sets, Set, Sets, card, cards, Card, Cards, board, boards, Board, Boards, select, selects, Select, Selects, identify, identifies, Identify, Identifies, find, finds, Find, Finds, choose, chooses, Choose, Chooses, match, matches, Match, Matches, matching, Matching, valid, Valid, invalid, Invalid, attribute, attributes, Attribute, Attributes, animal, animals, Animal, Animals, biome, biomes, Biome, Biomes, same, Same, different, Different, all, All, triple, triples, Triple, Triples, triad, triads, Triad, Triads, combination, combinations, Combination, Combinations, group, groups, Group, Groups, game, games, Game, Games, gameset, gamesets, Gameset, Gamesets, rule, rules, Rule, Rules, constraint, constraints, Constraint, Constraints, value, values, Value, Values \\[2pt]
Operator Function & arithmetic, Arithmetic, operator, operators, Operator, Operators, function, functions, Function, Functions, calculation, calculations, Calculation, Calculations, compute, Compute, calculate, Calculate, evaluate, Evaluate, solve, Solve, multiplication, multiplications, Multiplication, Multiplications, multiply, Multiply, product, products, Product, Products, addition, additions, Addition, Additions, add, Add, sum, sums, Sum, Sums, integer, integers, Integer, Integers, input, inputs, Input, Inputs, result, results, Result, Results, expression, expressions, Expression, Expressions, composition, compositions, Composition, Compositions \\[2pt]
Noun Class Agreement & grammar, Grammar, grammars, Grammars, grammaticality, Grammaticality, grammatical, Grammatical, ungrammatical, Ungrammatical, syntax, Syntax, sentence, Sentence, sentences, Sentences, phrase, Phrase, phrases, Phrases, noun, Noun, nouns, Nouns, determiner, Determiner, determiners, Determiners, class, Class, classes, Classes, agreement, Agreement, match, Match, matches, Matches, matching, Matching, mismatch, Mismatch, mismatches, Mismatches, template, Template, templates, Templates, validation, Validation, classification, Classification, decision, Decision, decisions, Decisions, binary, Binary, wellformedness, Wellformedness, acceptability, Acceptability \\[2pt]
Tapatan & board, boards, Board, Boards, game, games, Game, Games, grid, grids, Grid, Grids, cell, cells, Cell, Cells, point, points, Point, Points, coordinate, coordinates, Coordinate, Coordinates, piece, pieces, Piece, Pieces, player, players, Player, Players, state, states, State, States, outcome, outcomes, Outcome, Outcomes, result, results, Result, Results, winner, winners, Winner, Winners, win, wins, Win, Wins, winning, Winning, line, lines, Line, Lines, alignment, alignments, Alignment, Alignments, sequence, sequences, Sequence, Sequences, straight, Straight, horizontal, Horizontal, vertical, Vertical, diagonal, diagonals, Diagonal, Diagonals, orthogonal, Orthogonal, neighbor, neighbors, Neighbor, Neighbors, adjacency, adjacencies, Adjacency, Adjacencies, adjacent, Adjacent, connection, connections, Connection, Connections, connectivity, Connectivity, parity, parities, Parity, Parities, even, Even, empty, Empty, validity, Validity, valid, Valid, classification, classifications, Classification, Classifications, classify, Classify, determine, Determine, detect, Detect, evaluate, Evaluate, identify, Identify, label, labels, Label, Labels \\[2pt]
Lexical Category Inference & validation, Validation, validate, Validate, verification, Verification, assessment, Assessment, decision, Decision, classification, Classification, classify, Classify, categorization, Categorization, categorize, Categorize, category, Category, categories, Categories, membership, Membership, member, Member, members, Members, list, List, lists, Lists, item, Item, items, Items, word, Word, words, Words, order, Order, ordering, Ordering, sequence, Sequence, position, Position, positions, Positions, alignment, Alignment, matching, Matching, match, Match, criteria, Criteria, criterion, Criterion, rule, Rule, rules, Rules, constraint, Constraint, constraints, Constraints, definition, Definition, definitions, Definitions, correctness, Correctness, incorrectness, Incorrectness, satisfaction, Satisfaction, validity, Validity \\[2pt]
\bottomrule
\end{tabular}
\end{footnotesize}

\end{center}


\clearpage
\section*{Decision log}\label{app:decisions}
\addcontentsline{toc}{section}{Decision log}
\begin{description}
\item[2026-10-27] Scope: test the lexical-task-head method of \citet{yang2026shared} on complex tasks and compare the per-head logit lens with Head-JLens. New title.
\item[2026-10-27] Tasks: the 14 simple tasks shared by \citet{yang2026shared} and \citet{davidson2025different}; the five tasks of \citet{fu2026rules}.
\item[2026-10-27] No main condition: rule-based and example-based prompts are treated alike (replacing the earlier choice of rules as main and examples as secondary).
\item[2026-10-27] Method: lexical task heads, with Logit-head and J-head readouts. Function vectors, function-specific heads and the hidden-state difference between prompts are not used for now.
\item[2026-10-27] At least six rule wordings per complex task and difficulty level.
\item[2026-10-27] Term lists are task-specific, not item-specific; one per task, or per difficulty level where the rules change.
\item[2026-10-27] Answer format as in the source studies.
\item[2026-10-27] Model: Llama-3.1-8B-Instruct, the main model of \citet{yang2026shared} and the only one with released Head-JLens matrices.
\item[2026-10-27] Rule wordings and term lists for the complex tasks are generated with GPT-5.5 through ChatGPT and checked by hand.
\item[2026-10-27] Accuracy is computed with each source paper's evaluation.
\item[2026-10-27] Prompts on which the model does not answer straight away are left out; other ways of handling them are considered after E0.
\item[2026-10-27] After an early E0 finding (no direct answers on Set Game and Operator Function), these tasks are run in reasoning mode (Sec.~\ref{sec:reasoning}).
\item[2026-10-27] Tapatan is run in reasoning mode as well: in E0 none of its 300 items is answered straight away, under any prompt.
\item[2026-10-27] Complex tasks: the easy level of each task only.
\item[2026-10-27] Rule wordings and term lists are generated with GPT-5.5 through the Azure OpenAI API, replacing ChatGPT; six wordings per original rule prompt. Both were checked by hand and approved.
\item[2026-10-27] Complex tasks are decoded greedily, and use 300 test items per task drawn with the sampler and seed of \citet{fu2026rules} (all 600 for Lexical Category Inference), so that the causal check can leave the patched items out of the mean.
\item[2026-10-27] Reasoning mode: 100 items per task, style and wording; the causal check patches 10 items with $\alpha\in\{0,1,4,10\}$.
\end{description}

\end{document}
